\subsection{Actual original rational-unit rays and genuine integer branch balancing}

The complete four-node mathematical plan and both independent mathematical
reviews preceded code. The preceding original multiset and ANY-ring actual
principal-open unit modules are reused unchanged. Every function below is
the literal native regular-section unit logarithm, compared on its whole
nonzero locus with the ORIGINAL polynomial log difference. No profile,
local-slope, branch, positivity, comparison or harmonicity certificate is input.
Use an ACTUAL algebraically closed field K, ANY original valuation, positive
radius, and ORIGINAL nonzero polynomials p,q; no completeness or nontrivial
valuation is assumed.

\begin{theorem}[Actual Gauss rays derive all integer slopes of the same original unit logarithm]
\label{analytic:actual-original-rational-unit-branch-slopes}
\leanok
\lean{EnvelopeProperty.Berkovich.RelativeSpectrum.centeredGauss_val_X_sub_C_eq_max,
EnvelopeProperty.Berkovich.RelativeSpectrum.polynomialRealLog,
EnvelopeProperty.Berkovich.RelativeSpectrum.val_eq_leadingCoeff_mul_root_prod,
EnvelopeProperty.Berkovich.RelativeSpectrum.linear_val_ne_zero_of_polynomial_ne_zero,
EnvelopeProperty.Berkovich.RelativeSpectrum.polynomialRealLog_eq_root_sum,
EnvelopeProperty.Berkovich.RelativeSpectrum.polynomialRealLog_centeredGauss,
EnvelopeProperty.Berkovich.RelativeSpectrum.exists_centered_root_radius,
EnvelopeProperty.Berkovich.RelativeSpectrum.polynomialRealLog_centeredGauss_sub,
EnvelopeProperty.Berkovich.RelativeSpectrum.rationalRealLog,
EnvelopeProperty.Berkovich.RelativeSpectrum.sectionLogAbs_ratio_eq_rationalRealLog,
EnvelopeProperty.Berkovich.RelativeSpectrum.centeredDiscPoint,
EnvelopeProperty.Berkovich.RelativeSpectrum.centeredDiscPoint_mem_branch,
EnvelopeProperty.Berkovich.RelativeSpectrum.centeredDiscPoint_endpoint,
EnvelopeProperty.Berkovich.RelativeSpectrum.exists_rational_inward_slope,
EnvelopeProperty.Berkovich.RelativeSpectrum.exists_outer_root_radius,
EnvelopeProperty.Berkovich.RelativeSpectrum.polynomialRealLog_gauss,
EnvelopeProperty.Berkovich.RelativeSpectrum.polynomialRealLog_gauss_sub,
EnvelopeProperty.Berkovich.RelativeSpectrum.exists_rational_outward_slope,
EnvelopeProperty.Berkovich.RelativeSpectrum.not_root_near_of_branch_not_mem,
EnvelopeProperty.Berkovich.RelativeSpectrum.val_linear_eq_gauss_on_branch,
EnvelopeProperty.Berkovich.RelativeSpectrum.val_polynomial_eq_gauss_on_unlisted_branch,
EnvelopeProperty.Berkovich.RelativeSpectrum.rationalRealLog_eq_gauss_on_unlisted_branch,
EnvelopeProperty.Berkovich.RelativeSpectrum.rationalUnitLog,
EnvelopeProperty.Berkovich.RelativeSpectrum.rationalUnitLog_eq,
EnvelopeProperty.Berkovich.RelativeSpectrum.exists_rationalUnitLog_inward_slope,
EnvelopeProperty.Berkovich.RelativeSpectrum.exists_rationalUnitLog_outward_slope,
EnvelopeProperty.Berkovich.RelativeSpectrum.rationalUnitLog_eq_on_unlisted_branch}
\uses{analytic:actual-original-root-branch-partition,analytic:actual-original-rational-section-unit,spectrum:centered-gauss-whole-ray,spectrum:centered-gauss-generic-support,analytic:actual-disc-polynomial-root-dominance,analytic:actual-polynomial-gauss-inflation}
For EACH original $a$ with $v(a)\le r$, construct a real original radius
$0<s_a<r$ such that, for EVERY actual $s_a\le t\le r$, the ORIGINAL
centered Gauss point $\eta_{a,t}$ lies in the genuine unit locus and
\[
 h(\eta_{a,t})=h(\eta_r)
   +\bigl(m_q(B_a)-m_p(B_a)\bigr)(\log r-\log t).
\]
For $t<r$, these are genuine whole points of the ACTUAL branch $B_a$;
at $t=r$ the whole original point equals $\eta_r$.
Thus the derived inward integer slope, with length oriented AWAY from
$\eta_r$, is $\sigma(B_a)=m_q(B_a)-m_p(B_a)$.

Construct ALSO $R>r$ such that for EVERY actual $r\le t\le R$, the ORIGINAL
outer Gauss ray satisfies
\[
 h(\eta_t)=h(\eta_r)+\sigma_{\rm out}(\log t-\log r),
 \qquad
 \sigma_{\rm out}=\#\{b\in\operatorname{roots}(p):v(b)\le r\}
                 -\#\{b\in\operatorname{roots}(q):v(b)\le r\}.
\]
EVERY actual unlisted branch $B_a\notin S$ has the SAME original function
$h(x)=h(\eta_r)$ at ALL its whole original points; in particular its slope
is zero. No source profile, positivity, unit, radius-neighborhood or slope
certificate is supplied.
\end{theorem}
\begin{proof}
The genuine entire split factorization of each original polynomial and
actual multiplicativity compute its value at the original centered Gauss
point from EVERY actual root. The linear-factor identity is exactly
$|(T-b)(\eta_{a,t})|=\max(t,v(a-b))$, obtained from the original whole
translation and actual Gauss coefficient formula.

There are only finitely many roots of $p,q$ with $v(a-b)<r$. Their actual
strict distances have a finite maximum below $r$; choose $s_a$ strictly
between that maximum (and zero) and $r$. Every matched linear factor has
value $t$ for $s_a\le t\le r$. Every unmatched linear factor has distance
at least $r$, hence constant value $v(a-b)$ throughout that interval.
At $t=r$, the actual original centered Gauss point equals $\eta_r$ since
$v(a)\le r$. Genuine factorization, positivity and the real log product
identity therefore give the displayed formula with TRUE original counts.
The actual centered coordinate value is $t$ and original coordinate value
is at most $\max(t,v(a))\le r$; this proves exact branch membership for
$t<r$. Positive radii make every nonzero polynomial value nonzero, so all
these whole original points lie in the actual unit locus.

For the outer ray, finitely many ORIGINAL roots have $v(b)>r$. Choose
$R>r$ no larger than every such value, inserting $r+1$ if the list is empty.
For $r\le t\le R$, each root inside the closed disc has factor value $t$,
and every outside root contributes its constant original value. Apply
whole factorization and real logs to derive the second formula. No claim
that an outer ray is itself an affinoid domain is needed.

Finally suppose $B_a\notin S$. Every original root has $v(a-b)\ge r$;
otherwise its actual branch would equal $B_a$ and occur in $S$. For ANY
whole $x\in B_a$, the actual strict bound $|(T-a)(x)|<r\le v(a-b)$
and unequal-term valuation identity force
$|(T-b)(x)|=v(a-b)=\max(r,v(b))$, the true original Gauss factor value.
All original polynomial factors therefore have their same nonzero Gauss
values. This proves that the ENTIRE original unit-log function on this
branch is constant, not merely constant along the displayed ray.
\end{proof}

\begin{theorem}[A true finite integer balancing identity at the actual original Gauss point]
\label{analytic:actual-original-rational-unit-gauss-balancing}
\leanok
\lean{EnvelopeProperty.Berkovich.RelativeSpectrum.exists_rationalUnitLog_balanced_rays}
\uses{analytic:actual-original-root-branch-partition,analytic:actual-original-rational-unit-branch-slopes}
For the SAME original regular unit, genuine branches and derived integer
slopes, the ACTUAL finite supported family satisfies
\[
 \sigma_{\rm out}+\sum_{B\in S}\sigma(B)=0.
\]
Every other ACTUAL inward connected-component branch has derived slope zero.
This balances the SAME original function, all original root multiplicities,
the genuine inward branches and the specified true outer Gauss direction.
It is not an assumed harmonicity record or a conclusion from radial
convexity alone.
\end{theorem}
\begin{proof}
The genuine finite root partition proves
$\sum_{B\in S}m_p(B)=n_p$ and $\sum_{B\in S}m_q(B)=n_q$,
where $n_p,n_q$ are the ACTUAL closed-disc root counts from the outer slope.
The proved inward orientation gives $\sigma(B)=m_q(B)-m_p(B)$;
the true outer orientation gives $\sigma_{\rm out}=n_p-n_q$.
Integer finite-sum subtraction and these genuine partition identities give
$(n_p-n_q)+(n_q-n_p)=0$. The preceding WHOLE-branch constancy proves the
zero slopes on every unlisted branch. No desired balancing identity, local
profile or synthetic branch partition was supplied as an input.
\end{proof}

\begin{remark}[Exact potential-theory scope]
This constructs actual integer slopes/balancing of original rational-unit
logarithms in genuine A1 branch geometry. It does not construct arbitrary
analytic-unit factorization, a semistable curve skeleton, residue-curve
weights, the general curve harmonic sheaf, actual relative analytic boundary
or Thuillier subharmonic comparison/finite descent. The outer term is the
specified original Gauss ray; classification of every branch of the WHOLE
projective line remains separate. The existing same-section and whole-point
identities keep this producer attached to the unchanged original functions.
\end{remark}
